\chapter{Graph}

\section{Fundamentals}
	\kactlimport{BellmanFord.h}
	\kactlimport{FloydWarshall.h}
	\kactlimport{TopoSort.h}

\section{Network flow}
	\kactlimport{PushRelabel.h}
	\kactlimport{MinCostMaxFlow.h}
	\kactlimport{EdmondsKarp.h}
	% \kactlimport{Dinic.h}
	\kactlimport{MinCut.h}
	\kactlimport{GlobalMinCut.h}
	\kactlimport{GomoryHu.h}

\section{Matching}
	\subsection{Matching Basics}
	\begin{itemize}
		\item To achieve maximum matching in a bipartite graph, keep altering augmenting path until no more exist. 
		\item \textbf{Hall's Theorem}: A bipartite graph $G$ with parts $X$ and $Y$ has a matching that covers all vertices in $X$ iff for every subset $S \subseteq X$, the number of neighbors of $S$ is at least as large as the size of $S$.
		\item \textbf{Kőnig's Theorem}: In any bipartite graph, the number of edges in a maximum matching equals the number of vertices in a minimum vertex cover, and it is complementary to the maximum independent set.
		\item \textbf{Node-disjoint path cover (for DAG)}: We can find a minimum node-disjoint path cover by constructing a matching graph where each node of the original graph is represented by two nodes: a left node and a right node. There is an edge from a left node to a right node if there is such an edge in the original graph. Minimum node-disjoint path cover = number of nodes - maximum matching in the matching graph.
		\item \textbf{General path cover (for DAG)} A minimum general path cover can be found almost like a minimum node-disjoint path cover. It suffices to add some new edges to the matching graph so there is an edge a → b always when there is a path from a to b in the original graph (possibly through several edges).
		\item \textbf{Dilworth's theorem}: In a DAG, the size of the minimum general path cover equals the size of the maximum antichain (set of nodes such that no node is reachable from another).
	\end{itemize}
	\kactlimport{HopcroftKarp.h}
	\kactlimport{DFSMatching.h}
	\kactlimport{MinimumVertexCover.h}
	\kactlimport{WeightedMatching.h}
	\kactlimport{GeneralMatching.h}

\section{DFS algorithms}
	\kactlimport{SCC.h}
	\kactlimport{BiconnectedComponents.h}
	\kactlimport{2sat.h}
	\kactlimport{EulerWalk.h}

\section{Coloring}
	\kactlimport{EdgeColoring.h}

\section{Heuristics}
	\kactlimport{MaximalCliques.h}
	\kactlimport{MaximumClique.h}
	\kactlimport{MaximumIndependentSet.h}

\section{Trees}
	\kactlimport{BinaryLifting.h}
	\kactlimport{LCA.h}
	\kactlimport{CompressTree.h}
	\kactlimport{HLD.h}
	\kactlimport{LinkCutTree.h}
	\kactlimport{DirectedMST.h}

\section{Math}
	\subsection{Number of Spanning Trees}
		% I.e. matrix-tree theorem.
		% Source: https://en.wikipedia.org/wiki/Kirchhoff%27s_theorem
		% Test: stress-tests/graph/matrix-tree.cpp
		Create an $N\times N$ matrix \texttt{mat}, and for each edge $a \rightarrow b \in G$, do
		\texttt{mat[a][b]--, mat[b][b]++} (and \texttt{mat[b][a]--, mat[a][a]++} if $G$ is undirected).
		Remove the $i$th row and column and take the determinant; this yields the number of directed spanning trees rooted at $i$
		(if $G$ is undirected, remove any row/column).

	\subsection{Erdős–Gallai theorem}
		% Source: https://en.wikipedia.org/wiki/Erd%C5%91s%E2%80%93Gallai_theorem
		% Test: stress-tests/graph/erdos-gallai.cpp
		A simple graph with node degrees $d_1 \ge \dots \ge d_n$ exists iff $d_1 + \dots + d_n$ is even and for every $k = 1\dots n$,
		\[ \sum _{i=1}^{k}d_{i}\leq k(k-1)+\sum _{i=k+1}^{n}\min(d_{i},k). \]
